(1) 并联元件伸长相同，弹簧力为 \(kx\)，粘性力为 \(b\dot x\)。忽略惯性时力平衡 \(b\dot x+kx=F_0\)，初值 \(x(0)=0\)，时间常数 \(\tau=b/k\)；这里恒力加载并不意味着初始瞬间出现有限位移跳跃。
(2) 一阶线性方程解为 \(x(t)=(F_0/k)(1-e^{-t/\tau})\)。弹簧承担 \(F_s=kx=F_0(1-e^{-t/\tau})\)，阻尼器承担 \(F_d=b\dot x=F_0e^{-t/\tau}\)，任意时刻 \(F_s+F_d=F_0\)，加载开始力全由粘性支路承担，久后全由弹簧承担。
(3) 最终位移 \(x_\infty=F_0/k\)，外力总功 \(W_{\rm ext}=\int_0^\infty F_0\dot x\,dt=F_0^2/k\)。储存弹性能 \(U_s=kx_\infty^2/2=F_0^2/(2k)\)，粘性损耗 \(\int_0^\infty b\dot x^2dt=(F_0^2/b)\int_0^\infty e^{-2t/\tau}dt=F_0^2/(2k)\)，故 \(W_{\rm ext}=U_s+Q_{\rm diss}\)。
(4) 完全蠕变后撤力、将撤力时刻记作新时间 \(t=0\)，自由回复满足 \(b\dot x+kx=0\) 且 \(x(0)=F_0/k\)，从而 \(x(t)=F_0e^{-t/\tau}/k\)；原有弹性能最终全部转化为追加热耗散 \(F_0^2/(2k)\)。\(b\to0\) 时 \(\tau\to0\)，出现理想化瞬时形变；\(b\to\infty\) 时任意固定时间内几乎不动，不能先取无限时间再把两极限交换。
